% Labels and cross-references refer to paper/paper.tex.
% This is a separate explanatory fragment, not an additional paper input.
% Formal-development details relocated from the paper.
\section{The Lean formalization}\label{sec:lean}
The declarations below connect the certificate checks with the mathematical theorem.

\subsection{The statement}
The definitions use Mathlib's unordered pairs \texttt{Sym2}; a coloring of
$K_n$ assigns a color to every pair, and the diagonal pairs are never
used.
\begin{Verbatim}
abbrev Col (n : ℕ) := Sym2 (Fin n) → Fin 3

def HasMonoPath (n m : ℕ) (X : Col n) : Prop :=
  ∃ (c : Fin 3) (p : Fin m → Fin n), Function.Injective p ∧
    ∀ (i : ℕ) (h : i + 1 < m), X s(p ⟨i, by omega⟩, p ⟨i + 1, h⟩) = c

noncomputable def ramsey3Path (m : ℕ) : ℕ :=
  sInf {n | ∀ X : Col n, HasMonoPath n m X}

theorem K18_has_mono_P10 (X : Sym2 (Fin 18) → Fin 3) :
    ∃ (c : Fin 3) (p : Fin 10 → Fin 18), Function.Injective p ∧
      ∀ (i : ℕ) (h : i + 1 < 10), X s(p ⟨i, by omega⟩, p ⟨i + 1, h⟩) = c

theorem K17_no_mono_P10 : ¬ HasMonoPath 17 10 k17Col

theorem R3_P10_eq_18 : ramsey3Path 10 = 18
\end{Verbatim}
A path is an injective map $p$ from $\{0,\dots,m-1\}$ to the vertices whose
consecutive pairs all have color $c$, which is the usual (not necessarily
induced) path on $m$ vertices. \texttt{ramsey3Path m} is the least $n$ for
which every coloring of $K_n$ has a monochromatic $P_m$; the proof of
\texttt{R3\_P10\_eq\_18} shows that $18$ belongs to that set and that no
smaller $n$ does, since a good coloring of $K_{17}$ restricts to one of any
smaller complete graph.

\subsection{What is proved}
Lean defines a closed-form encoder \texttt{encode n k}, which produces
the formula of Section~\ref{sec:encoding}, and proves for all $n$ and $k$
that it is sound for the colorings: a coloring with the fixed path
that satisfies every generator inequality, with the auxiliaries set to
their meaning, satisfies every clause (\texttt{encode\_sat}); a good
coloring satisfies every path clause (\texttt{paths\_ok}). It proves
Lemma~\ref{lem:fix} (\texttt{relabel}), that every generator preserves the
fixed path and goodness (\texttt{gens\_fix}, \texttt{act\_noMono}),
Lemma~\ref{lem:lex} (\texttt{orbit\_min}), the rung argument
(\texttt{rung}), Lemma~\ref{lem:cover} (\texttt{unsat\_of\_cubes}), the
chain of rungs, and the lower bound (Section~\ref{sec:lower}). The facts
about formulas are stated with the DIMACS semantics: a literal is a nonzero
integer, an assignment is a predicate on variables, and
\[
\texttt{UnsatWithPaths n m F}
\]
says that $F$ together with \emph{every} path clause for $P_m$ in $K_n$
has no model.

The proof of \texttt{K18\_has\_mono\_P10} uses nine hypotheses of this
kind: one for each of the seven rungs, one for the cover formula
\texttt{encode 18 9 ++ finalCubes.map negCube}, and one for every cube $c$
of \texttt{finalCubes}, \texttt{encode 18 9 ++ cubeUnits c} (one
universally quantified hypothesis over the list, so $15{,}964$ instances
in all). The list \texttt{finalCubes} is parsed in Lean from the
accompanying cube lists, and the cover formula and the cube hypothesis
use this same list.
The conditional theorems \texttt{K18\_P10\_of\_certs} and
\texttt{R3\_P10\_eq\_18\_of\_certs} take the nine facts as hypotheses and
depend only on Lean's standard axioms \texttt{propext},
\texttt{Classical.choice} and \texttt{Quot.sound}, as does
\texttt{K17\_no\_mono\_P10}.

\subsection{How the certificates enter}\label{sec:certs}
Each instance is proved by a Boolean check
\texttt{checkUnsat n m F paths proof} and its soundness theorem
\[
\texttt{checkUnsat\_sound :\ checkUnsat n m F paths proof = true →
UnsatWithPaths n m F},
\]
proved in Lean. The check confirms that every literal of $F$ is nonzero;
that each entry of \texttt{paths} (a color and $m$ distinct vertices below
$n$) yields a path clause, which Lean rebuilds from the path itself; and
that Lean's verified LRAT checker (\texttt{Std.Tactic.BVDecide.LRAT}, with
its soundness theorem) accepts \texttt{proof} as a refutation of $F$
together with those path clauses. The bridge between the two semantics is
a proved lemma (\texttt{toStd\_eval}): under the nonzero-literal condition,
an assignment satisfying the formula in the DIMACS semantics above gives an
assignment satisfying its translation into the checker's clause type, so
the checker's soundness theorem yields \texttt{UnsatWithPaths n m F}, the
same predicate that the rung and cover arguments consume. Only the path
clauses that the trimmed proof uses are included. Nothing outside Lean decides what counts as a
path clause, and the formula checked is a Lean term, not a file: the base
formula used in the cube checks is parsed from the accompanying $K_{18}$
formula and proved equal to \texttt{encode 18 9} by a comparison that is
itself natively evaluated.

Each certificate is stored as a compact string (an LRAT proof with
delta-coded hints and implicit clause identifiers). It is decoded by total
functions written in Lean. The decoder is outside the logical trusted
base: whatever it produces, acceptance implies unsatisfiability by
\texttt{checkUnsat\_sound}. The check is a closed Boolean term,
and the kernel accepts its value through \texttt{Lean.ofReduceBool}, which
evaluates it with Lean's compiled code. The $15{,}956$ cube facts are
checked in $162$ modules, each covering a contiguous range of indices; a
final module joins the ranges into \texttt{K18.all}.

\subsection{Axioms and native evaluation}
The axiom report of both main theorems is
\begin{Verbatim}
'R3P10.K18_has_mono_P10' depends on axioms:
  [propext, Classical.choice, Lean.ofReduceBool, Lean.trustCompiler, Quot.sound]
'R3P10.R3_P10_eq_18' depends on axioms:
  [propext, Classical.choice, Lean.ofReduceBool, Lean.trustCompiler, Quot.sound]
\end{Verbatim}
The two non-standard axioms come from the natively evaluated certificate
checks (Lean v4.32.2, pinned in the accompanying project). The trusted base
therefore has three parts: Lean's kernel, which checks all the mathematics,
including the lower bound and the reduction from colorings to formulas;
Lean's compiler and interpreter, whose evaluation of the closed Boolean
certificate checks the kernel accepts without reducing them itself
(\texttt{Lean.ofReduceBool}, \texttt{Lean.trustCompiler}); and the
compiled modules a build imports, which Lean does not recheck. The evaluated code
is the verified LRAT checker, the decoder and the clause and path builders;
by \texttt{checkUnsat\_sound}, only the faithful execution of this code is
trusted, not its correctness. Lean does not recheck imported compiled
modules, so a reader who wants the whole proof checked on one machine
rebuilds the $162$ certificate modules (see the accompanying checking record). The SAT
solver, the proof converter, the external checkers and the scripts that
produced the certificates are not part of the trusted base of the Lean
theorem.



\section{Additional formal conventions}
In Lean the coloring is written out edge by edge, and the argument is a
general lemma (for all $n$ and $m$: if each color has a labeling of the
vertices that is constant along edges of that color and whose classes have
fewer than $m$ vertices, there is no monochromatic $P_m$) together with two
finite facts decided by the kernel's own reduction (\texttt{decide
+kernel}), without native evaluation.


The following two conventions refer to the fixed-path and lex-leader
lemmas in the paper, respectively.
\begin{quote}
Here ``good'' means having no monochromatic $P_m$, for the same $m$
as in the preceding encoding. The following lemma fixes only path edges.

In the following lemma, goodness still refers to $P_m$ and $\Gamma$ is
the preceding list of path-preserving vertex and color permutations.
\end{quote}

\section{Further certificate and implementation details}
The formal statement agrees with Theorem~\ref{thm:main} as follows.
In namespace \texttt{R3P10}, the theorem without certificate hypotheses is
\texttt{R3\_P10\_eq\_18 : ramsey3Path 10 = 18}.
Here \texttt{Col n} is \texttt{Sym2 (Fin n) → Fin 3}: colors on unordered
vertex pairs. Diagonal values are unused and arbitrary; colors $0,1,2$
replace $1,2,3$. The predicate \texttt{HasMonoPath n m X} means that there
exist a color and an injective map from $m$ vertices to $n$ vertices whose
consecutive pairs have that color. It requires no condition on other
pairs, so it means our non-induced $P_m$, with $m$ vertices, not edges.
The definition \texttt{ramsey3Path m} is the infimum in the natural numbers
of the $n$ for which every \texttt{Col n} satisfies that predicate.
The proof establishes membership at $18$ and a countercoloring at $17$;
restriction of that coloring excludes every smaller $n$. Thus this
infimum is precisely the Ramsey number in Theorem~\ref{thm:main}, with
no remaining certificate hypotheses.

Lean proves the lower bound, symmetry reductions, encoding soundness and
assembly of refutations. Its verified LRAT checker checks propositional refutations in which each
step cites earlier clauses by number. The final theorem's axioms are
\texttt{propext}, \texttt{Classical.choice}, \texttt{Quot.sound},
\texttt{Lean.ofReduceBool} and \texttt{Lean.trustCompiler}.
The last two support evaluation by compiled code of certificate checks,
cube-list lengths and equality of the parsed and defined base formulas. The proof trusts Lean's kernel,
compiler and imported compiled modules; rebuilding the certificate modules
avoids trusting previously built copies. The SAT solver, converter and
external checkers are outside this trusted base. Literature attributions,
historical timings and reports of separate checker executions are outside
Lean. The final section of this note gives the formal declarations; the accompanying
checking notes give the certificate-import details.


The recorded complete computation checked all cases twice,
in 249 and 301 CPU-hours. The second used six independent batches of
12 CPU cores; its longest batch took 7.5 hours. The accompanying
completion record gives the phase timings and software versions. The
archived certificates \cite{Gra26} allow the certificate checks to be
repeated without repeating the search. The default Lean build proves
the conditional theorem; building the final theorem requires all the
certificate modules. The README gives the full and sample-check commands
and the requirements.



\subsection{Path clauses generated during the search}\label{sec:prop}
The solver is CaDiCaL~\cite{BFF+24} with an external propagator attached
through the IPASIR-UP interface~\cite{FNP+23,FNP+24}. The propagator
follows the assignment of the color variables. Whenever an edge receives a
color $c$, it searches for paths on $m$ vertices through that edge, in the
graph of edges not yet excluded from color $c$, all of whose edges except
at most one are already known to have color $c$. If all edges are known,
it reports the path clause as a conflict; if exactly one is unknown, it
sets that edge to ``not $c$'' with the path clause as the reason. When the
solver reaches a complete assignment, the propagator searches it
exhaustively for a monochromatic $P_m$ and, if it finds one, rejects the
assignment with that path clause. Every clause handed to the solver is a
path clause and is written to a log.

The refuted formula is therefore the base formula plus the logged path
clauses (plus a cube's unit clauses), a subset of the formula with
\emph{all} path clauses. A refutation of the subset refutes the full
formula; nothing about the completeness of the propagator is needed.
Conversely, a satisfying assignment is accepted only after the exhaustive
search, and every coloring the method reports is rechecked by a separate
program.




\subsection{The cubes}
The
cubes were produced by the solver itself, with the propagator attached.
Walking the edges in edge order, at each edge it assumes each color in
turn and runs unit propagation (including the path propagator); colors
leading to a conflict are dropped, an edge with a single surviving color
is fixed without counting as a split, and the walk stops after $10$ edges
with at least two surviving colors. This produced $15{,}316$ cubes of
$11$ to $20$ literals, and $18$ path clauses learned during generation.

Sixteen of these cubes, each with $11$ literals, are listed in
\path{data/k18/split/parents.txt}. Each was split once more in the same way, by four further split
edges, into $41$ children, $656$ in all ($144$ more path clauses were
learned during this split). The final list consists of the $15{,}300$
remaining first-level cubes and the $656$ children: $15{,}956$ cubes.

\subsection{The cover formula}\label{sec:cover}
Cubes whose propagation fails are dropped during generation, so the list is
not complete by construction. Its completeness is proved instead by one
more certificate: the \emph{cover formula} of Section~\ref{sec:cnc}, which consists of the base
formula, the $162$ path clauses learned during generation, and the
negation of every final cube (one clause per cube), $30{,}790$ clauses in
all, is unsatisfiable.

With $F$ the $K_{18}$ base formula, Lemma~\ref{lem:cover} and the
$15{,}957$ certificates give the unsatisfiability of the top rung, and the
successive bounds then give $\R(P_{10})\le18$. In Lean this is the theorem combining
the cover certificate with the per-cube certificates
(\texttt{unsat\_of\_cubes}); the per-cube facts enter through
\texttt{cubes\_of\_range}, which turns a theorem
\texttt{K18.all :\ ∀ i, 0 ≤ i → i < 15956 → CubeOK i} over the list's
indices into the statement for every member of the list, given that the
list has length $15{,}956$.

\section{Formal declarations}\label{note:declarations}
The declarations below give the definitions used by the formal theorem.
The coloring \texttt{k17Col} is the coloring of Section~\ref{sec:lower}.

\begin{Verbatim}
abbrev Col (n : ℕ) := Sym2 (Fin n) → Fin 3

def HasMonoPath (n m : ℕ) (X : Col n) : Prop :=
  ∃ (c : Fin 3) (p : Fin m → Fin n), Function.Injective p ∧
    ∀ (i : ℕ) (h : i + 1 < m), X s(p ⟨i, by omega⟩, p ⟨i + 1, h⟩) = c

noncomputable def ramsey3Path (m : ℕ) : ℕ :=
  sInf {n | ∀ X : Col n, HasMonoPath n m X}

theorem K18_has_mono_P10 (X : Sym2 (Fin 18) → Fin 3) :
    ∃ (c : Fin 3) (p : Fin 10 → Fin 18), Function.Injective p ∧
      ∀ (i : ℕ) (h : i + 1 < 10), X s(p ⟨i, by omega⟩, p ⟨i + 1, h⟩) = c

theorem K17_no_mono_P10 : ¬ HasMonoPath 17 10 k17Col

theorem R3_P10_eq_18 : ramsey3Path 10 = 18
\end{Verbatim}

\section{Complete verification statement and historical measurements}
The formal theorem is
\texttt{R3P10.R3\_P10\_eq\_18 : ramsey3Path 10 = 18}.
Its colorings assign one of the colors $0,1,2$ to every unordered pair of
vertices, including unused diagonal pairs. Relabeling the colors as
$1,2,3$ gives the convention of Theorem~\ref{thm:main}.
A formal monochromatic path is an injective map from $\{0,\dots,m-1\}$
whose consecutive pairs all have one color, with no condition on the
other pairs. Thus it is the same non-induced path on $m$ vertices used
here. The formal Ramsey number is the natural-number infimum of the
orders at which every coloring contains such a path. The proof shows
that this set contains $18$ and excludes every smaller order by
restricting the countercoloring on $17$ vertices. The final theorem has
no certificate hypotheses; the finite coloring in its lower bound is
exactly the construction of Section~\ref{sec:lower}.

Lean proves the reductions, the lower bound and the soundness of the
certificate checker. The certificate checks, the
cube-list lengths and the equality of the parsed and defined formulas
are evaluated by compiled code. The permitted axioms are
\texttt{propext}, \texttt{Classical.choice}, \texttt{Quot.sound},
\texttt{Lean.ofReduceBool} and \texttt{Lean.trustCompiler}.
Accordingly, the proof depends on Lean's kernel, compiler and imported
compiled modules. The SAT solver and proof converter produce certificates
whose correctness is checked, rather than assumed; the separate checker
executions, literature attributions and timing records are outside Lean.

The default Lean build proves the conditional theorem; building the
final theorem requires all the certificate modules.

The complete computation was checked twice, in 249 and 301 CPU-hours.
The second used six independent batches of 12 CPU cores; its longest
batch took 7.5 hours. Source code, the completed-computation record and
checking instructions accompany the paper. The archived
certificates~\cite{Gra26} permit checking without repeating the search.


\subsection{Previously quoted source passages}
For this construction, Gyárfás, Ruszinkó, Sárközy and Szemerédi
\cite[author manuscript, July 9, 2012, Section~1.1, paragraph following
Theorem~1]{GRSS07} write: ``One type comes from a `blow-up' of a
factorization of $K_4$.'' The odd and even part sizes are specified there;
see also Dybizbański, Dzido and Radziszowski~\cite[arXiv:1207.5851v3,
Section~2, paragraph following Conjecture~2]{DDR16}.

The following excerpts and locators are retained from the earlier presentation; the paper now cites their mathematical content in prose. The full texts of Böving et al.\ \cite{BBC+25} and Kirchweger, Manrique and Szeider \cite{KMS26} were unavailable in the presentation review, so those quotations have not been independently reauthenticated here.

The advance is the upper bound $\R(P_{10})\le18$.
Gyárfás, Ruszinkó, Sárközy and Szemerédi proved
$\R(P_n)=2n-2+(n\bmod2)$ for all sufficiently large $n$
\cite[Theorem~1]{GRSS07}, with its corrigendum.
Dybizbański, Dzido and Radziszowski proved $\R(P_8)=14$
(Theorem~10) and $\R(P_9)\le17$ (Theorem~8), which with their Section~2
lower bound gives $\R(P_9)=17$ \cite{DDR16}.
Radziszowski's survey \cite[Section~6.4, p.~65]{Rad26}
records the formula through $n=9$ and $P_{10}$ as the first open case.
Our $K_{17}$ lower bound is the even-case blow-up in
\cite[Section~1.1, after Theorem~1]{GRSS07}, also given in
\cite[Section~2]{DDR16}.
In their Conjecture~2, Dybizbański, Dzido and Radziszowski attribute
the suggestion to Faudree and Schelp (1975), quoting ``they feel''.

This paper settles $n=10$ of Conjecture~\ref{conj:FS}; the general
conjecture remains open. The difficult step is obtaining a tractable case
split without losing a coloring: fixing a guaranteed path and imposing
symmetry constraints must be jointly sound, and discarded branches leave
a coverage obligation. A lexicographic minimum and a separate refutation
excluding all assignments outside the selected cases discharge these
obligations. Dybizbański, Dzido and
Radziszowski's closing remarks identify the difficulty:
``the complexity of case analysis appears to be quite harder''
\cite[Section~5]{DDR16}. Referring to their attempts at $P_{10}$ and
$P_{11}$, they also state, ``We did not complete these proofs''
\cite[Section~5]{DDR16}.

The method combines cube-and-conquer with symmetry breaking and verified
certificate checking. Heule, Kullmann, Wieringa and Biere's
cube-and-conquer \cite[Section~3]{HKWB12} ``partitions a formula into many
subformulas''. For certificate import we use the verified ``LRAT UNSAT
proof checker'' of Böving et al.\ \cite[introduction and Section~4.4]{BBC+25}.
LRAT certificates are propositional refutations whose steps refer to
earlier clauses by number.
Szeider's streaming tool \texttt{lrat-catcher} ``turns a certificate into a
Lean theorem'' \cite[abstract]{Sze26}; our import uses the same checker
with its own compact encoding. Kirchweger, Manrique and Szeider describe
``a unified, machine-checked verification pipeline'' for graph generation
\cite[abstract]{KMS26}.

Computational precedents include Heule, Kullmann and Marek's
Boolean Pythagorean triples proof \cite[Theorem~1]{HKM16}, Codish, Frank,
Itzhakov and Miller's $R(4,3,3)=30$ \cite[abstract]{CFIM16}, and Heule and
Scheucher's empty-hexagon theorem \cite{HS24}, with Subercaseaux et al.'s
formal verification \cite{SNG+24}. In related Ramsey computations, Wesley
proves $R(C_3,C_6,C_6)=R(C_5,C_6,C_6)=15$
\cite[Theorems~3 and~4]{Wes25}; Galindez's draft reports
``$R(C_{10},C_{10},C_{10})\le26$'' \cite[Theorem~1]{Gal26}, using
cube-and-conquer and Kirchweger--Szeider's ``SAT modulo symmetries''
\cite[title]{KS24}; it explicitly reports no DRAT/LRAT proof certificate.
Shao, Jiang, Xu and Wang provide a
``reproduction package: $R(C_4,C_6,C_8)=12$'' \cite[title]{SJXW26}, and
Leo-Y-Zhang provides \textit{CycleRamsey} \cite{Zha26}.
Zhang's tree-bound corollary $R_k(T)\le k(n-2)+2$ for trees on $n\ge2$
vertices is \texttt{multicolourRamsey\_le\_tight} in
\texttt{Erdos557/Tight.lean} \cite{ZhJ26}. It uses an AI-found
Erdős--Sós proof packaged by T.~Adamczewski and gives $\R(P_{10})\le26$.

Earlier path and cycle literature includes Dzido \cite{Dzi05},
Dzido--Fidytek \cite{DF09}, Omidi--Raeisi \cite{OR11},
Maherani et al.\ \cite[Theorem~1.1]{MORS15}, Benevides--Skokan
\cite{BS09}, Figaj--Łuczak \cite{FL07} and Letzter \cite{Let22}.
For small cases, see Irving, Yang--Rowlinson and Yang
\cite{Irv74,YR94,YR92,Yan95}, and Radziszowski's survey
\cite[p.~65, Section~6.4.1(c)]{Rad26}.
The extremal background is due to Erdős--Gallai \cite{EG59},
Faudree--Schelp \cite{FS75}, whose theorem is recalled by Dybizbański,
Dzido and Radziszowski \cite[Theorem~4]{DDR16}. See also Kopylov's ``Maximal paths and cycles in a graph'' \cite[title]{Kop77}.
Allen et al.\ study paths and cycles with few color changes
\cite[title]{ABC+26}.
We use de~Moura and Ullrich's Lean~4 \cite{MU21}, the mathlib Community's
library \cite{Mathlib20}, Wetzler, Heule and Hunt's \texttt{drat-trim}
\cite{WHH14} for conversion to Cruz-Filipe et al.'s LRAT format
\cite{CHH+17}, and Tan, Heule and Myreen's \texttt{cake\_lpr}
\cite{THM21} for the separate checks.

\section{Source locators for the paper}
The following pinpoints accompany the ordinary citations in the paper. Retained published-v5 locator strings reproduce that version's wording: its description of a README as ``pinned'' has no immutable upstream commit in the available evidence.
\begin{itemize}
\item \cite{DDR16}: Preprint locators: Conjecture~2, Theorems~4, 8 and~10; Section~2 (lower construction), Section~4 (the $P_9$ bound) and Section~5 (closing remarks on $P_{10}$ and $P_{11}$).
\item \cite{GRSS07}: Author-manuscript locators: Theorem~1 and Section~1.1, paragraph following Theorem~1 (the odd and even blow-up constructions).
\item \cite{Rad26}: Locator: p.~65, Sections~6.4 and~6.4.1(c).
\item \cite{HKWB12}: Method locator: Section~3.
\item \cite{BBC+25}: Checker locators: introduction and Section~4.4.
\item \cite{Sze26}: Streaming-import description: abstract.
\item \cite{KMS26}: Graph-generation description: abstract.
\item \cite{Gal26}: Draft note, August 18, 2026, Theorem~1 and Section~5.3; the limitations explicitly report no DRAT/LRAT certificate.
\item \cite{CGLR96}: Encoding locator: Section~3, Proposition~3.1 and the following auxiliary-variable clauses.
\item \cite{ZhJ26}: Tree-bound corollary: \texttt{Erdos557/Tight.lean}, \texttt{multicolourRamsey\_le\_tight}, commit \texttt{9d26a79d9ad04ff605167d022e512984a76ea3fa}; imported-proof attribution: the README at that commit.
\item \cite{CFIM16}: abstract (the value $R(4,3,3)=30$).
\item \cite{GRSS07}: retained published-v5 pinpoints: author manuscript, July 9, 2012, Theorem~1.
\item \cite{DDR16}: retained published-v5 pinpoints: arXiv:1207.5851v3, Theorems~10 and~8 with Section~2; arXiv:1207.5851v3, Section~2; arXiv:1207.5851v3, Theorem~4; arXiv:1207.5851v3, Section~5; arXiv:1207.5851v3, Section~4.
\item \cite{Rad26}: retained published-v5 pinpoints: revision~18, p.~65, Section~6.4.1(c); revision~18, p.~65, Sections~6.4 and~6.4.1(c).
\item \cite{HKWB12}: retained published-v5 pinpoints: author manuscript, Section~3, Overview.
\item \cite{BBC+25}: retained published-v5 pinpoints: version of record, introduction, contributions, and
Section~4.4.
\item \cite{Sze26}: retained published-v5 pinpoints: arXiv:2607.00815v2, abstract.
\item \cite{KMS26}: retained published-v5 pinpoints: IJCAR 2026, abstract.
\item \cite{HKM16}: retained published-v5 pinpoints: author preprint, Theorem~1.
\item \cite{CFIM16}: retained published-v5 pinpoints: author preprint, abstract.
\item \cite{HS24}: retained published-v5 pinpoints: TACAS 2024, title.
\item \cite{SNG+24}: retained published-v5 pinpoints: ITP 2024, title.
\item \cite{Wes25}: retained published-v5 pinpoints: arXiv:2509.03784v1.
\item \cite{Gal26}: retained published-v5 pinpoints: Zenodo record title.
\item \cite{KS24}: retained published-v5 pinpoints: TOCL 2024, title.
\item \cite{SJXW26}: retained published-v5 pinpoints: Zenodo record title.
\item \cite{Zha26}: retained published-v5 pinpoints: repository title.
\item \cite{ZhJ26}: retained published-v5 pinpoints: pinned README, Ramsey corollary.
\item \cite{ABC+26}: retained published-v5 pinpoints: arXiv:2607.03243v1, title.
\item \cite{MORS15}: retained published-v5 pinpoints: arXiv:1207.3771v1, Theorem~1.1.
\item \cite{CGLR96}: retained published-v5 pinpoints: technical report TR-1996-012, Section~3,
Proposition~3.1 and the following auxiliary-variable clauses.
\end{itemize}
